\chapter{Transformations du plan}
\textit{Translater une figure, la retourner devant un miroir, l'agrandir, la faire
tourner : toutes ces opérations sont des \emph{transformations du plan}. Chacune
déplace les points selon une règle précise, mais respecte certaines propriétés
(distances, alignement, parallélisme, angles). Ce chapitre recense les
transformations au programme, leurs propriétés et leurs expressions analytiques —
des outils que tu retrouveras en géométrie, en analyse et plus tard avec les nombres
complexes.}

\section{Transformation du plan : vocabulaire}

\begin{definition}[Transformation du plan]
Une \textbf{transformation du plan} est une application $f$ qui, à tout point $M$ du
plan, associe un unique point $M'=f(M)$, appelé \textbf{image} de $M$, et telle que
tout point du plan possède \emph{exactement un} antécédent. On dit que $M'$ est
l'image de $M$ et que $M$ est l'antécédent de $M'$.
\end{definition}

\begin{definition}[Point invariant]
Un point $M$ est dit \textbf{invariant} (ou \textbf{fixe}) par la transformation $f$
si $f(M)=M$.
\end{definition}

\begin{aretenir}[Propriétés que peut conserver une transformation]
Selon la transformation, certaines de ces propriétés sont conservées :
\begin{itemize}
  \item les \textbf{distances} ($MN = M'N'$) : on parle alors d'\textbf{isométrie} ;
  \item l'\textbf{alignement} de trois points ;
  \item le \textbf{parallélisme} de deux droites ;
  \item les \textbf{mesures d'angles} ;
  \item le \textbf{barycentre} (l'image du barycentre est le barycentre des images,
        avec les mêmes coefficients).
\end{itemize}
\end{aretenir}

\begin{remarque}
Toutes les transformations de ce chapitre conservent l'alignement, le parallélisme
et le barycentre. Translation, symétries et rotation conservent \emph{aussi} les
distances (isométries). L'homothétie, elle, \emph{multiplie} les distances par
$|k|$ : ce n'est pas une isométrie (sauf si $k=\pm 1$).
\end{remarque}

\section{La translation}

\begin{definition}[Translation de vecteur $\vec{u}$]
Soit $\vec{u}$ un vecteur du plan. La \textbf{translation de vecteur $\vec{u}$},
notée $t_{\vec{u}}$, est la transformation qui à tout point $M$ associe l'unique
point $M'$ tel que
\[ \overrightarrow{MM'}=\vec{u}. \]
\end{definition}

\begin{propriete}[Propriétés de la translation]
Soit $t_{\vec u}$ la translation de vecteur $\vec u$.
\begin{itemize}
  \item Si $\vec u\neq\vec 0$, $t_{\vec u}$ n'a \textbf{aucun point invariant} ;
        si $\vec u=\vec 0$, c'est l'identité (tout point est invariant).
  \item C'est une \textbf{isométrie} : pour tous $M,N$ d'images $M',N'$,
        $\ M'N'=MN$. En effet $\overrightarrow{M'N'}=\overrightarrow{MN}$.
  \item Elle conserve l'\textbf{alignement}, le \textbf{parallélisme}, les
        \textbf{angles} et le \textbf{barycentre}.
\end{itemize}
\end{propriete}

\begin{theoreme}[Conservation du vecteur]
Pour tous points $M$ et $N$ d'images $M'$ et $N'$ par $t_{\vec u}$, on a
$\overrightarrow{M'N'}=\overrightarrow{MN}$.
\end{theoreme}

\begin{remarque}
\textit{Démonstration.} Par définition $\overrightarrow{MM'}=\vec u$ et
$\overrightarrow{NN'}=\vec u$. Avec la relation de Chasles,
$\overrightarrow{M'N'}=\overrightarrow{M'M}+\overrightarrow{MN}+\overrightarrow{NN'}
=-\vec u+\overrightarrow{MN}+\vec u=\overrightarrow{MN}$. D'où $M'N'=MN$.
\end{remarque}

\begin{propriete}[Expression analytique]
Le plan est muni d'un repère $(O,\vec{\imath},\vec{\jmath})$ et $\vec u(a\,;\,b)$.
La translation $t_{\vec u}$ transforme $M(x\,;\,y)$ en $M'(x'\,;\,y')$ avec
\[ \begin{cases} x'=x+a\\[2pt] y'=y+b. \end{cases} \]
\end{propriete}

\begin{exemple}
Avec $\vec u(3\,;\,-1)$, l'image de $A(2\,;\,5)$ est $A'(2+3\,;\,5-1)=A'(5\,;\,4)$.
\end{exemple}

\section{La symétrie centrale}

\begin{definition}[Symétrie de centre $O$]
Soit $O$ un point du plan. La \textbf{symétrie centrale de centre $O$}, notée $S_O$,
associe à tout point $M$ l'unique point $M'$ tel que $O$ soit le \textbf{milieu} du
segment $[MM']$, c'est-à-dire $\overrightarrow{OM'}=-\overrightarrow{OM}$.
\end{definition}

\begin{propriete}[Propriétés de la symétrie centrale]
\begin{itemize}
  \item Le \textbf{seul point invariant} est le centre $O$.
  \item C'est une \textbf{isométrie} : $M'N'=MN$. En effet
        $\overrightarrow{M'N'}=-\overrightarrow{MN}$.
  \item Elle conserve l'alignement, le parallélisme, les angles et le barycentre.
  \item $S_O$ est sa \textbf{propre réciproque} : $S_O\circ S_O=\mathrm{id}$.
  \item La symétrie centrale est l'\textbf{homothétie de centre $O$ et de rapport
        $-1$}.
\end{itemize}
\end{propriete}

\begin{propriete}[Expression analytique]
Si $O(x_0\,;\,y_0)$, la symétrie $S_O$ transforme $M(x\,;\,y)$ en
$M'(x'\,;\,y')$ avec
\[ \begin{cases} x'=2x_0-x\\[2pt] y'=2y_0-y. \end{cases} \]
En particulier, si le centre est l'origine $O$, alors $x'=-x$ et $y'=-y$.
\end{propriete}

\begin{exemple}
Image de $A(4\,;\,1)$ par la symétrie de centre $O(1\,;\,2)$ :
$x'=2(1)-4=-2$, $y'=2(2)-1=3$, donc $A'(-2\,;\,3)$.
\end{exemple}

\section{La symétrie orthogonale (axiale)}

\begin{definition}[Symétrie d'axe $(D)$]
Soit $(D)$ une droite. La \textbf{symétrie orthogonale d'axe $(D)$} (ou
\textbf{symétrie axiale}), notée $S_{(D)}$, associe à tout point $M$ le point $M'$
tel que $(D)$ soit la \textbf{médiatrice} du segment $[MM']$. Si $M\in(D)$, alors
$M'=M$.
\end{definition}

\begin{propriete}[Propriétés de la symétrie axiale]
\begin{itemize}
  \item Les \textbf{points invariants} sont exactement les points de l'axe $(D)$.
  \item C'est une \textbf{isométrie} : elle conserve les distances.
  \item Elle conserve l'alignement, le parallélisme, les mesures d'angles
        (mais \textbf{inverse leur orientation}) et le barycentre.
  \item $S_{(D)}$ est sa propre réciproque : $S_{(D)}\circ S_{(D)}=\mathrm{id}$.
\end{itemize}
\end{propriete}

\begin{propriete}[Expressions analytiques usuelles]
Dans un repère orthonormé, $M(x\,;\,y)$ a pour image $M'(x'\,;\,y')$ :
\begin{center}\renewcommand{\arraystretch}{1.3}
\begin{tabular}{@{}l l@{}}
\toprule
\textbf{Axe de symétrie} & \textbf{Image $M'(x'\,;\,y')$} \\
\midrule
axe des abscisses $(Ox)$ & $x'=x,\quad y'=-y$ \\
axe des ordonnées $(Oy)$ & $x'=-x,\quad y'=y$ \\
première bissectrice $y=x$ & $x'=y,\quad y'=x$ \\
\bottomrule
\end{tabular}
\end{center}
\end{propriete}

\begin{remarque}
La composée de deux symétries axiales d'axes parallèles est une \textbf{translation} ;
la composée de deux symétries axiales d'axes sécants en $O$ est une
\textbf{rotation} de centre $O$.
\end{remarque}

\section{L'homothétie}

\begin{definition}[Homothétie de centre $\Omega$ et de rapport $k$]
Soit $\Omega$ un point et $k$ un réel \textbf{non nul}. L'\textbf{homothétie de
centre $\Omega$ et de rapport $k$}, notée $h(\Omega,k)$, associe à tout point $M$
l'unique point $M'$ tel que
\[ \overrightarrow{\Omega M'}=k\,\overrightarrow{\Omega M}. \]
\end{definition}

\begin{propriete}[Propriétés de l'homothétie]
Soit $h=h(\Omega,k)$, $k\neq 0$.
\begin{itemize}
  \item Le centre $\Omega$ est invariant. Si $k\neq 1$, c'est le \textbf{seul}
        point invariant.
  \item Pour tous $M,N$ d'images $M',N'$ : $\overrightarrow{M'N'}=k\,\overrightarrow{MN}$,
        donc $M'N'=|k|\,MN$ (les longueurs sont \textbf{multipliées par $|k|$}).
  \item Elle conserve l'\textbf{alignement}, le \textbf{parallélisme}
        (une droite a pour image une droite \emph{parallèle}), les \textbf{angles}
        et le \textbf{barycentre}.
  \item Cas particuliers : $k=1$ donne l'identité ; $k=-1$ donne la \textbf{symétrie
        centrale} de centre $\Omega$.
\end{itemize}
\end{propriete}

\begin{theoreme}[Conservation du parallélisme]
Pour tous points $M,N$ d'images $M',N'$ par $h(\Omega,k)$, on a
$\overrightarrow{M'N'}=k\,\overrightarrow{MN}$.
\end{theoreme}

\begin{remarque}
\textit{Démonstration.} Par définition $\overrightarrow{\Omega M'}=k\overrightarrow{\Omega M}$
et $\overrightarrow{\Omega N'}=k\overrightarrow{\Omega N}$. Alors, avec Chasles,
$\overrightarrow{M'N'}=\overrightarrow{\Omega N'}-\overrightarrow{\Omega M'}
=k(\overrightarrow{\Omega N}-\overrightarrow{\Omega M})=k\,\overrightarrow{MN}$.
Les vecteurs $\overrightarrow{M'N'}$ et $\overrightarrow{MN}$ sont colinéaires :
$(M'N')\parallel(MN)$, et $M'N'=|k|\,MN$.
\end{remarque}

\begin{illus}
\begin{tikzpicture}[scale=1,>={Stealth[]}]
  \coordinate (Om) at (0,0);
  \coordinate (A) at (1.2,0.6);
  \coordinate (B) at (2.2,0.3);
  \coordinate (C) at (1.7,1.5);
  % image par k=1.8
  \coordinate (Ap) at (2.16,1.08);
  \coordinate (Bp) at (3.96,0.54);
  \coordinate (Cp) at (3.06,2.7);
  \filldraw[onbink] (Om) circle (1.6pt) node[below left,onbink]{$\Omega$};
  \draw[thick,onbmuted] (A)--(B)--(C)--cycle;
  \filldraw[onbo] (A) circle (1.3pt) node[below,onbink]{$A$};
  \filldraw[onbo] (B) circle (1.3pt) node[right,onbink]{$B$};
  \filldraw[onbo] (C) circle (1.3pt) node[above,onbink]{$C$};
  \draw[thick,onbo] (Ap)--(Bp)--(Cp)--cycle;
  \filldraw[onbo] (Ap) circle (1.3pt) node[below,onbink]{$A'$};
  \filldraw[onbo] (Bp) circle (1.3pt) node[right,onbink]{$B'$};
  \filldraw[onbo] (Cp) circle (1.3pt) node[above,onbink]{$C'$};
  \draw[->,onbblue,dashed] (Om)--(Cp);
  \draw[->,onbline] (Om)--(A);
\end{tikzpicture}
\fig{Figure — image du triangle $ABC$ par l'homothétie de centre $\Omega$ et de rapport $k=1{,}8$ : la figure est agrandie, chaque côté reste parallèle à son image.}
\end{illus}

\begin{propriete}[Expression analytique]
Si $\Omega(x_0\,;\,y_0)$, l'homothétie $h(\Omega,k)$ transforme $M(x\,;\,y)$ en
$M'(x'\,;\,y')$ avec
\[ \begin{cases} x'=x_0+k(x-x_0)\\[2pt] y'=y_0+k(y-y_0). \end{cases} \]
\end{propriete}

\begin{exemple}
Image de $A(5\,;\,2)$ par $h(\Omega,2)$ avec $\Omega(1\,;\,1)$ :
$x'=1+2(5-1)=9$, $y'=1+2(2-1)=3$, donc $A'(9\,;\,3)$.
\end{exemple}

\section{La rotation}

\begin{definition}[Rotation de centre $\Omega$ et d'angle $\theta$]
Soit $\Omega$ un point et $\theta$ un angle orienté. La \textbf{rotation de centre
$\Omega$ et d'angle $\theta$}, notée $r(\Omega,\theta)$, associe à tout point $M$ :
\begin{itemize}
  \item $\Omega$ lui-même si $M=\Omega$ ;
  \item sinon, l'unique point $M'$ tel que $\Omega M'=\Omega M$ et
        $\big(\overrightarrow{\Omega M},\overrightarrow{\Omega M'}\big)=\theta$.
\end{itemize}
\end{definition}

\begin{propriete}[Propriétés de la rotation]
\begin{itemize}
  \item Le \textbf{centre $\Omega$} est le seul point invariant (si $\theta$ n'est
        pas un multiple de $2\pi$).
  \item C'est une \textbf{isométrie} : elle conserve les distances.
  \item Elle conserve l'alignement, le parallélisme, les \textbf{mesures et
        l'orientation des angles} et le barycentre.
  \item Cas particuliers : $\theta=0$ donne l'identité ; $\theta=\pi$ donne la
        \textbf{symétrie centrale} de centre $\Omega$.
  \item La réciproque de $r(\Omega,\theta)$ est $r(\Omega,-\theta)$.
\end{itemize}
\end{propriete}

\begin{illus}
\begin{tikzpicture}[scale=1,>={Stealth[]}]
  \coordinate (Om) at (0,0);
  \coordinate (M) at (2.4,0);
  \coordinate (Mp) at (60:2.4);
  \filldraw[onbink] (Om) circle (1.6pt) node[below left,onbink]{$\Omega$};
  \draw[onbline,dashed] (Om) circle (2.4);
  \draw[->,onbmuted] (Om)--(M) node[right,onbink]{$M$};
  \draw[->,onbo,thick] (Om)--(Mp) node[above right,onbink]{$M'$};
  \filldraw[onbo] (M) circle (1.3pt);
  \filldraw[onbo] (Mp) circle (1.3pt);
  \draw[onbblue] (1,0) arc (0:60:1);
  \node[onbblue] at (40:1.35) {$\theta$};
\end{tikzpicture}
\fig{Figure — rotation de centre $\Omega$ et d'angle $\theta$ : $M'$ est sur le cercle de centre $\Omega$ passant par $M$, avec $\big(\protect\overrightarrow{\Omega M},\protect\overrightarrow{\Omega M'}\big)=\theta$.}
\end{illus}

\begin{remarque}
Au programme de Première C, on retient surtout la \emph{définition géométrique} et
les propriétés de la rotation ; l'expression analytique générale (avec $\cos\theta$
et $\sin\theta$) sera systématisée en Terminale, notamment grâce aux nombres
complexes.
\end{remarque}

\section{Image d'une droite, d'un cercle}

\begin{theoreme}[Image d'une droite]
Par une translation, une symétrie (centrale ou axiale), une homothétie ou une
rotation, l'image d'une \textbf{droite} est une \textbf{droite}. Plus précisément :
\begin{itemize}
  \item par une translation, l'image de $(D)$ est une droite \textbf{parallèle} à $(D)$ ;
  \item par une homothétie, l'image de $(D)$ est une droite \textbf{parallèle} à $(D)$ ;
  \item par une symétrie axiale ou une rotation, l'image est une droite (en général
        non parallèle, sauf cas particuliers).
\end{itemize}
\end{theoreme}

\begin{theoreme}[Image d'un cercle]
L'image d'un \textbf{cercle} $\mathcal{C}$ de centre $I$ et de rayon $R$ est un cercle
$\mathcal{C}'$ de centre $I'$ (image de $I$) et de rayon $R'$ :
\begin{itemize}
  \item par une \textbf{isométrie} (translation, symétrie, rotation) : $R'=R$ ;
  \item par une \textbf{homothétie} de rapport $k$ : $R'=|k|\,R$.
\end{itemize}
\end{theoreme}

\begin{exemple}
Soit $\mathcal{C}$ le cercle de centre $I(1\,;\,0)$ et de rayon $2$, et $h$
l'homothétie de centre $O$ et de rapport $3$. L'image de $I$ est $I'(3\,;\,0)$ et le
rayon devient $|3|\times 2=6$ : $\mathcal{C}'$ a pour centre $I'(3\,;\,0)$ et rayon $6$.
\end{exemple}

\section{Composées simples}

\begin{propriete}[Quelques composées à connaître]
\begin{itemize}
  \item $t_{\vec u}\circ t_{\vec v}=t_{\vec u+\vec v}$ : la composée de deux
        translations est la translation de vecteur $\vec u+\vec v$.
  \item $h(\Omega,k')\circ h(\Omega,k)=h(\Omega,kk')$ : deux homothéties de
        \emph{même centre} se composent en multipliant les rapports.
  \item $S_O\circ S_O=\mathrm{id}$ et $S_{(D)}\circ S_{(D)}=\mathrm{id}$ (les
        symétries sont \textbf{involutives}).
  \item $r(\Omega,\theta')\circ r(\Omega,\theta)=r(\Omega,\theta+\theta')$ : deux
        rotations de même centre s'additionnent en angle.
\end{itemize}
\end{propriete}

\begin{exemple}
La composée de la translation de vecteur $\vec u(1\,;\,2)$ suivie de celle de vecteur
$\vec v(3\,;\,-5)$ est la translation de vecteur $\vec u+\vec v=(4\,;\,-3)$.
\end{exemple}

\begin{remarque}
Attention : la composition de transformations n'est pas toujours commutative
($f\circ g\neq g\circ f$ en général). Elle l'est cependant pour deux translations,
ou pour deux homothéties (ou rotations) de \emph{même centre}.
\end{remarque}

% ════════════════════════════════════════════════════════════
\section{L'essentiel à retenir}

\begin{synthese}
\textbf{Définitions vectorielles.}
$t_{\vec u}:\overrightarrow{MM'}=\vec u$\ ;\quad
$S_O:\overrightarrow{OM'}=-\overrightarrow{OM}$\ ;\quad
$h(\Omega,k):\overrightarrow{\Omega M'}=k\,\overrightarrow{\Omega M}$.

\smallskip
\textbf{Tableau récapitulatif.}
\begin{center}\renewcommand{\arraystretch}{1.35}
\color{white}
\begin{tabular}{@{}l c c c@{}}
\toprule
\textbf{Transformation} & \textbf{Distances} & \textbf{Invariant(s)} & \textbf{Effet} \\
\midrule
Translation $t_{\vec u}$        & conservées      & aucun (si $\vec u\neq\vec 0$) & isométrie \\
Symétrie centrale $S_O$         & conservées      & le centre $O$                 & isométrie \\
Symétrie axiale $S_{(D)}$       & conservées      & les points de $(D)$           & isométrie \\
Rotation $r(\Omega,\theta)$     & conservées      & le centre $\Omega$            & isométrie \\
Homothétie $h(\Omega,k)$        & $\times|k|$     & le centre $\Omega$            & $R'=|k|R$ \\
\bottomrule
\end{tabular}
\end{center}

\smallskip
\textbf{Toutes} conservent alignement, parallélisme, angles et barycentre.

\smallskip
\textbf{Expressions analytiques.}
$t_{\vec u}:x'=x+a,\ y'=y+b$\ ;\quad
$S_O$ (centre $O$) : $x'=-x,\ y'=-y$\ ;\quad
$S_{(Ox)}:y'=-y$, $S_{(Oy)}:x'=-x$, $S_{y=x}:(x',y')=(y,x)$\ ;\quad
$h(\Omega,k):x'=x_0+k(x-x_0),\ y'=y_0+k(y-y_0)$.

\smallskip
\textbf{Images.} Droite $\to$ droite (parallèle pour $t_{\vec u}$ et $h$) ;
cercle $\to$ cercle (rayon $R$, ou $|k|R$ pour $h$).

\smallskip
\textbf{Composées.} $t_{\vec u}\circ t_{\vec v}=t_{\vec u+\vec v}$ ;
$h(\Omega,k')\circ h(\Omega,k)=h(\Omega,kk')$ ;
$S\circ S=\mathrm{id}$ ;
$r(\Omega,\theta')\circ r(\Omega,\theta)=r(\Omega,\theta+\theta')$.
\end{synthese}

% ════════════════════════════════════════════════════════════
\section{Méthodes \& savoir-faire}

\methode{Construire l'image d'un point par une transformation}
Revenir à la \emph{définition vectorielle} : reporter $\vec u$ depuis $M$
(translation), prendre le symétrique par rapport à $O$ ou à $(D)$ (symétries), ou
placer $M'$ sur la demi-droite $[\Omega M)$ avec $\Omega M'=|k|\,\Omega M$ (homothétie,
en tenant compte du signe de $k$).
\begin{exemple}
Pour $h(\Omega,-2)$ : $M'$ est tel que $\overrightarrow{\Omega M'}=-2\,\overrightarrow{\Omega M}$,
donc $M'$ est de l'autre côté de $\Omega$, à une distance $\Omega M'=2\,\Omega M$.
\end{exemple}

\methode{Calculer l'image d'un point par l'expression analytique}
Identifier la transformation et ses paramètres ($\vec u$, centre, rapport),
appliquer directement les formules sur les coordonnées de $M(x\,;\,y)$.
\begin{exemple}
Image de $A(-1\,;\,4)$ par $h(\Omega,3)$, $\Omega(2\,;\,2)$ :
$x'=2+3(-1-2)=-7$, $y'=2+3(4-2)=8$, soit $A'(-7\,;\,8)$.
\end{exemple}

\methode{Reconnaître la nature d'une transformation par sa relation}
Comparer la relation donnée aux définitions : $\overrightarrow{MM'}$ constant
$\Rightarrow$ translation ; $\overrightarrow{\Omega M'}=k\,\overrightarrow{\Omega M}$
$\Rightarrow$ homothétie (symétrie centrale si $k=-1$) ; distances conservées avec un
seul point fixe et angle constant $\Rightarrow$ rotation.
\begin{exemple}
Si $\overrightarrow{AM'}=\tfrac{1}{2}\overrightarrow{AM}$ pour tout $M$, la
transformation est l'homothétie de centre $A$ et de rapport $\tfrac{1}{2}$ : elle
réduit les longueurs de moitié.
\end{exemple}

\methode{Déterminer l'image d'une droite ou d'un cercle}
Pour une droite, il suffit de l'image de \emph{deux} de ses points (ou d'un point et
de la direction). Pour un cercle, transformer le \emph{centre} puis ajuster le rayon
($R$ pour une isométrie, $|k|R$ pour une homothétie).
\begin{exemple}
Image du cercle de centre $I(0\,;\,3)$, rayon $4$, par la translation de vecteur
$\vec u(2\,;\,-1)$ : centre $I'(2\,;\,2)$, rayon inchangé $4$.
\end{exemple}

\methode{Composer deux transformations}
Appliquer la première transformation, puis la seconde au résultat ; pour les cas
usuels (deux translations, deux homothéties/rotations de même centre), utiliser
directement les règles de composition.
\begin{exemple}
$r(\Omega,\frac{\pi}{3})\circ r(\Omega,\frac{\pi}{6})=r(\Omega,\frac{\pi}{6}+\frac{\pi}{3})
=r(\Omega,\frac{\pi}{2})$ : une seule rotation d'un quart de tour.
\end{exemple}

\methode{Exploiter la conservation du barycentre}
Comme l'image d'un barycentre est le barycentre des images (mêmes coefficients), pour
transformer le milieu, l'isobarycentre ou un point particulier d'une figure, il
suffit de transformer les sommets puis de reformer le barycentre.
\begin{exemple}
Si $I$ est le milieu de $[AB]$, alors son image $I'$ par toute transformation du
chapitre est le milieu de $[A'B']$ : on transforme $A$ et $B$, et $I'$ s'en déduit.
\end{exemple}
